% Figures and text: IVF search with per-list query blocks and two-pass pruning (ann_ivf_ttaiss)
% Written for readers who know ANN/IVF search
% but not the Tenstorrent architecture: the algorithm first, the hardware in
% one paragraph. Requires \usetikzlibrary{decorations.pathreplacing} as well.
%
% Uses the thesis colours hostc, devc and the styles darr, sectag; the boxes are
% defined locally. Requires \usetikzlibrary{positioning,calc,arrows.meta} and
% amsmath. If the thesis styles are not available:
% \definecolor{hostc}{RGB}{214,228,240} \definecolor{devc}{RGB}{226,236,214}
% \tikzset{sectag/.style={font=\scriptsize,text=gray,inner sep=1pt},
%          darr/.style={-{Stealth[length=2mm]},thick}}

%------------------------------------------------------------------------------
% Figure 1: the algorithm, step by step
%------------------------------------------------------------------------------
\begin{figure}[htbp]
\centering
\resizebox{\textwidth}{!}{%
\begin{tikzpicture}[font=\small,
  stepbox/.style={draw,rounded corners=3pt,align=center,text width=2.75cm,minimum height=3.1cm,inner sep=4pt},
  cpu/.style={stepbox,fill=hostc!70},
  acc/.style={stepbox,fill=devc!80,thick},
  outp/.style={font=\scriptsize\itshape,align=center,text width=2.9cm}]
% steps
\node[acc] (s1) at (0,0) {\textbf{1. Coarse search}\\[3pt]
  score every query against all centroids; keep its \(N_{\mathrm{probe}}\) closest clusters, ranked};
\node[cpu] (s2) at (3.45,0) {\textbf{2. Group by cluster}\\[3pt]
  for each cluster, collect the queries that probe it; pack them into blocks of 32};
\node[acc] (s3) at (6.9,0) {\textbf{3. Pass 1}\\[3pt]
  each query's 8 closest clusters: exact scores, top-\(k\) per query};
\node[acc] (s4) at (10.35,0) {\textbf{4. Pass 2}\\[3pt]
  its other \(N_{\mathrm{probe}}-8\) clusters; skip every block that cannot beat the bound \(T_q\)};
\node[acc] (s5) at (13.8,0) {\textbf{5. Final top-\(k\)}\\[3pt]
  merge pass 1 and pass 2 into the top-\(k\) of each query};
\node[cpu] (s6) at (17.25,0) {\textbf{6. Results}\\[3pt]
  return the top-\(k\); optional FP32 re-ranking of the device top-32};
\draw[darr] (s1) -- (s2);
\draw[darr] (s2) -- (s3);
\draw[darr] (s3) -- (s4);
\draw[darr] (s4) -- (s5);
\draw[darr] (s5) -- (s6);
% outputs
\node[outp,anchor=north] at (s1.south) {\(N_{\mathrm{probe}}\) clusters per query, closest first};
\node[outp,anchor=north] at (s2.south) {blocks of \(\le 32\) queries per cluster};
\node[outp,anchor=north] at (s3.south) {provisional top-\(k\) per query; bound \(T_q\) = its \(k\)-th best score\(^{*}\)};
\node[outp,anchor=north] at (s4.south) {top-\(k\) of each (cluster, block) that was not skipped};
\node[outp,anchor=north] at (s5.south) {top-\(k\) per query (exact for its \(N_{\mathrm{probe}}\) clusters)};
\node[outp,anchor=north] at (s6.south) {\(k\) nearest neighbours per query};
% the bound
\draw[darr,dashed] (s3.north) to[bend left=25] node[above,font=\scriptsize]{bound \(T_q\) for every query} (s4.north);
% legend
\node[draw,fill=devc!80,thick,minimum width=0.5cm,minimum height=0.3cm] (la) at (0.4,-3.3) {};
\node[anchor=west,font=\scriptsize] at (la.east) {on the accelerator (64 cores)};
\node[draw,fill=hostc!70,minimum width=0.5cm,minimum height=0.3cm] (lc) at (4.4,-3.3) {};
\node[anchor=west,font=\scriptsize] at (lc.east) {on the host CPU (step 2 for pass 2 runs while the accelerator executes pass 1)};
\end{tikzpicture}}
\caption{One batch search. The coarse search selects the \(N_{\mathrm{probe}}\)
closest clusters of every query. Instead of scanning clusters query by query,
the queries are grouped by cluster, so each cluster is read once and multiplied
with dense blocks of the queries that need it. Pass 1 scans each query's 8
closest clusters and gives a provisional top-\(k\); its \(k\)-th best score
\(T_q\) is a lower bound on the final one, so in pass 2 any block whose scores
are all below the bounds of its queries cannot change the result and is
skipped. A final top-\(k\) merge combines both passes. \(^{*}\)The 32nd best score
when the host re-ranks the 32 best candidates.}
\label{fig:gf-algorithm}
\end{figure}

%------------------------------------------------------------------------------
% Figure 2: the two ideas
%------------------------------------------------------------------------------
\begin{figure}[htbp]
\centering
\resizebox{\textwidth}{!}{%
\begin{tikzpicture}[font=\small,
  cell/.style={draw,gray!60,minimum size=0.45cm,inner sep=0pt},
  mat/.style={draw,thick,align=center,font=\scriptsize}]
% ---------------- (a) cluster-major grouping
\node[anchor=west,font=\bfseries] at (-0.6,4.6) {(a) Group queries by cluster};
\foreach \j in {1,...,6} \node[font=\scriptsize] at (0.45*\j,4.05) {\(C_{\j}\)};
\foreach \i in {1,...,7} \node[font=\scriptsize] at (-0.1,3.75-0.45*\i) {\(q_{\i}\)};
\foreach \i in {1,...,7} \foreach \j in {1,...,6} \node[cell] at (0.45*\j,3.75-0.45*\i) {};
\fill[devc] (0.45*3-0.225,3.75-0.45*7-0.225) rectangle (0.45*3+0.225,3.75-0.45*1+0.225);
\foreach \i in {1,...,7} \node[cell] at (0.45*3,3.75-0.45*\i) {};
\foreach \i/\j in {1/1,1/3,1/5, 2/2,2/4,2/6, 3/1,3/2,3/5, 4/3,4/4,4/6, 5/1,5/2,5/4, 6/3,6/5,6/6, 7/2,7/3,7/4}
  \fill (0.45*\j,3.75-0.45*\i) circle (0.08);
\node[font=\scriptsize,align=center,text width=3.0cm] at (1.55,0.15)
  {dot = query probes cluster\\(\(N_{\mathrm{probe}}\) dots per row)};
\draw[darr] (3.05,2.4) -- node[above,font=\scriptsize,align=center]{column \(C_3\)} (4.25,2.4);
\node[mat,minimum width=1.1cm,minimum height=1.9cm,fill=devc!60] (qb) at (4.95,2.4)
  {queries\\probing\\\(C_3\)\\\(q_1,q_4,\)\\\(q_6,q_7,\ldots\)};
\node[font=\scriptsize] at (4.95,1.2) {block \(\le 32\)};
\node[font=\large] at (5.8,2.4) {\(\times\)};
\node[mat,minimum width=1.9cm,minimum height=1.1cm,fill=gray!10] (vb) at (7.0,2.4) {vectors of \(C_3\)\\(32 at a time)};
\node[font=\large] at (8.2,2.4) {\(=\)};
\node[mat,minimum width=1.3cm,minimum height=1.3cm,fill=white] (sb) at (9.1,2.4) {scores\\\(32\times32\)};
\draw[darr] (sb.south) -- (9.1,1.05);
\node[font=\scriptsize,align=center,text width=3.2cm] at (9.1,0.6) {update the top-\(k\)\\of each query in the block};
\node[font=\scriptsize,align=left,text width=5.6cm,anchor=north west] at (3.9,0.15)
  {Each cluster is read once per search and multiplied with all\\
   query blocks that need it; every row of a block is a query\\
   that actually probes the cluster.};
% ---------------- (b) two-pass pruning
\begin{scope}[xshift=11.6cm]
\node[anchor=west,font=\bfseries] at (-0.4,4.6) {(b) Two passes with a pruning bound (one query \(q\))};
\node[font=\scriptsize,anchor=west] at (-0.4,4.05) {clusters of \(q\), closest first (\(N_{\mathrm{probe}}=64\)):};
\foreach \j in {0,...,15} {
  \ifnum\j<2 \fill[devc] (0.38*\j,3.35) rectangle ++(0.38,0.42); \fi
  \draw[gray] (0.38*\j,3.35) rectangle ++(0.38,0.42);}
\draw[decorate,decoration={brace,amplitude=4pt,mirror}] (0,3.28) -- (0.76,3.28)
  node[midway,below=4pt,font=\scriptsize,align=center]{pass 1:\\8 closest};
\draw[decorate,decoration={brace,amplitude=4pt,mirror}] (0.8,3.28) -- (6.08,3.28)
  node[midway,below=4pt,font=\scriptsize]{pass 2: the other 56};
\node[font=\scriptsize,align=left,anchor=north west,text width=6.4cm] at (-0.4,2.35)
  {\textbf{Pass 1} gives a provisional top-\(k\) of \(q\). Its \(k\)-th best score
   \(T_q\) can only rise as more clusters are added, so a candidate scoring
   below \(T_q\) can never enter the final top-\(k\).};
\node[mat,minimum width=1.0cm,minimum height=1.0cm,fill=white] (b1) at (0.3,0.35) {block};
\node[font=\scriptsize,align=left,anchor=west,text width=5.0cm] at (0.95,0.35)
  {\textbf{Pass 2}, per block of 32 vectors \(\times\) 32 queries: if every
   score is below the bound of its query, skip the top-\(k\) update
   (at \(N_{\mathrm{probe}}=64\), most blocks are skipped).};
\end{scope}
\end{tikzpicture}}
\caption{The two ideas behind the method. (a)~Queries are grouped by cluster
instead of by arrival order: a block holds only queries that probe the
cluster, so the cluster's vectors are read once and almost every computed
score is needed. (b)~Each query is searched in two passes. The closest
clusters, which contain most true neighbours, are searched first; the
resulting bound lets the second pass skip blocks that cannot change the query's
top-\(k\). The final result equals that of a standard IVF search over the same
\(N_{\mathrm{probe}}\) clusters (up to equal scores).}
\label{fig:gf-ideas}
\end{figure}

%------------------------------------------------------------------------------
% Text
%------------------------------------------------------------------------------
\subsection{Algorithm}
\label{sec:gf-algorithm}

The index is a standard IVF index. The dataset vectors are partitioned into
\(N_{\mathrm{list}}\) clusters by k-means, and each cluster stores its vectors and
their IDs. Vectors are L2-normalised, so the score of a candidate is its inner
product with the query (cosine similarity). A search processes a batch of queries
(10{,}000 in our experiments) in five steps (Figure~\ref{fig:gf-algorithm}).

\begin{enumerate}
  \item \textbf{Coarse search.} Every query is scored against all centroids, and
    its \(N_{\mathrm{probe}}\) closest clusters are kept, ranked by centroid score.
  \item \textbf{Grouping by cluster.} The selection is inverted: for every
    cluster, the queries that probe it are collected and packed into blocks of
    32 queries.
  \item \textbf{Pass 1.} Each query is searched in its 8 closest clusters only
    (the first 8 of its ranked list; in general \(N_{\mathrm{probe}}/8\)). For
    every cluster and every block of its queries, the cluster's vectors are
    multiplied with the block, 32 vectors at a time, and each query keeps its
    top-\(k\) candidates. Merging a query's results over its 8 clusters gives a
    provisional top-\(k\). Its \(k\)-th best score is the query's bound \(T_q\)
    (the 32nd best when re-ranking follows, because then all 32 candidates
    passed to the host must be exact).
  \item \textbf{Pass 2.} The query's remaining \(N_{\mathrm{probe}}-8\) clusters
    are searched the same way, with one addition. For every block of 32 vectors
    \(\times\) 32 queries, the scores are first compared with the bounds of the
    block's queries. If no score exceeds its query's bound, the block cannot
    change any result and its top-\(k\) update is skipped.
  \item \textbf{Final top-\(k\).} For each query, the provisional top-\(k\) of pass 1
    and the per-cluster top-\(k\) lists of pass 2 are merged into its final
    top-\(k\). Optionally, the host re-scores the 32 best candidates in full FP32
    precision and keeps the best \(k\) (re-ranking).
\end{enumerate}

\paragraph{Why the result is exact.} Every query is compared only with the
vectors of its own \(N_{\mathrm{probe}}\) clusters, as in standard IVF. The bound
\(T_q\) is the \(k\)-th best score among a subset of the query's candidates, so
it can never exceed the \(k\)-th best score of the final result. A skipped
candidate scores at most \(T_q\), so it could only tie with the last result,
never displace it. The method therefore returns the same top-\(k\) as an IVF
search over the same clusters, up to candidates with equal scores. We verify
this against a reference search for every run.

\paragraph{Why it is fast.}
\begin{itemize}
  \item \emph{Each cluster is read once.} Grouping by cluster turns many
    query-by-query scans into a few dense matrix products: a cluster's vectors
    are loaded once per search and multiplied with all query blocks that need
    them.
  \item \emph{The blocks are dense.} A block holds only queries that probe the
    cluster, so about 90--99\% of the computed scores are used. A batch of 32
    consecutive queries would share few clusters, and most of its rows would be
    wasted.
  \item \emph{Most top-\(k\) updates are skipped.} Maintaining the top-\(k\) lists,
    not the distance computation, dominates the cost. A query's closest
    clusters hold most of its true neighbours, so the bound from pass 1 is
    already tight. At \(N_{\mathrm{probe}}=64\), only about 11\% of the pass-2
    blocks still need an update when the bound is the 32nd best score (as
    needed for re-ranking). An offline estimate gives about 3\% with the 10th
    best score.
\end{itemize}

\paragraph{Parallelism.} The accelerator, a Tenstorrent Wormhole, has 64 small
cores. Each core has its own memory and a unit that multiplies 32\(\times\)32
matrix tiles, which is why queries and vectors are handled in blocks of 32.
\begin{itemize}
  \item The coarse search and the final top-\(k\) are split over the cores by
    query: each core handles a contiguous range of 32-query groups.
  \item The passes are split by cluster. A work item is one cluster with up to
    32 of its query blocks, and the items are balanced over the 64 cores
    (largest first, onto the least-loaded core).
  \item Inside each core, loading the next vectors overlaps with computing on
    the current ones.
  \item Cores exchange results only through the accelerator's memory, between
    the steps.
  \item The host overlaps its own work with the accelerator: it prepares the
    grouping for pass 2 while pass 1 runs.
\end{itemize}
The steps run in order, because each step needs the previous one's output
(pass 2 needs the bounds of pass 1).
